A következő fejezet célja annak bemutatása, hogy a fejlesztés nem közvetlenül az implementációval indult, hanem egy előzetes terméktervezési folyamattal, amelynek központi eleme a Product Requirements Document (PRD).

A PRD alapján meghatározhatóvá vált a projekt kiinduló problémája, célközönsége, fő tervezett funkciói, valamint az első megvalósítható verzió, vagyis az MVP-hez (Minimum Viable Product) szükséges követelmények.

\section{A projekt kiinduló problémája}
Manapság az emberek filmnézési szokásai több platformra terjednek ki: streaming szolgáltatókra, mozikra, fizikai adathordozókra és televíziós tartalmakra. Ezek a megoldások rendkívül kézenfekvővé teszik a filmek fogyasztását, ugyanakkor egy új problémát is létrehoznak. Egy adott szolgáltatás, például a Netflix, képes nyilvántartani, hogy a felhasználó az adott platformon milyen filmeket nézett meg, vagy mit szeretne később megnézni, ez azonban nem terjed ki más szolgáltatásokra.

Emiatt a felhasználó filmnézési előzményei nehezen követhetővé válnak. Előfordulhat, hogy valaki nem emlékszik pontosan, látott-e már egy filmet, milyen értékelést adott volna rá, vagy mely filmeket szerette volna később megnézni. A probléma tehát nem kizárólag a filmek megtalálásáról szól, hanem egy személyes, platformfüggetlen filmnapló kialakításáról is.

A filmnézésnek ugyanakkor közösségi oldala is van. Gyakran kíváncsiak vagyunk arra, hogy ismerőseink milyen filmeket néztek meg az utóbbi időben, mit gondoltak róluk, vagy milyen filmeket ajánlanának. Ezek az ajánlások sokszor beszélgetésekben, csoportos üzenetekben vagy közösségi médiás bejegyzésekben jelennek meg, így könnyen elvesznek, és később nehezen kereshetők vissza. A legtöbb streaming platform elsősorban a saját katalógusára és az egyéni fogyasztásra koncentrál, ezért nem ad teljes választ arra, hogyan lehetne a személyes filmnaplót és a baráti ajánlásokat egységes rendszerben kezelni.

A WatchDB projekt kiinduló problémája ezért kettős. Egyrészt szükség van egy egyszerű, platformfüggetlen felületre, ahol a felhasználó nyilvántarthatja a megnézett filmeket és a később megtekintendő alkotásokat. Másrészt megjelenik az igény egy olyan közösségi rétegre is, amelyen keresztül a felhasználók követhetik egymás filmnézési aktivitását és ajánlásokat fedezhetnek fel. A dolgozatban bemutatott projekt ezt a problémát használja kiindulópontként, majd ebből vezeti le a termék koncepcióját, a követelményeket és az MVP-hez szükséges funkciókat.

\section{A PRD szerepe a fejlesztési folyamatban}

A Product Requirements Document (PRD) a projekt előkészítésének egyik legfontosabb dokumentuma. Szerepe, hogy a termékötletet strukturált formába rendezze, és összekapcsolja a felhasználói igényeket, az üzleti célokat, a funkcionális követelményeket és a fejlesztési korlátokat. A PRD így nem pusztán leíró dokumentum, hanem a fejlesztési folyamat egyik kiindulópontja: segít meghatározni, hogy milyen problémát kell megoldani, kinek készül a termék, milyen funkciók szükségesek az első verzióhoz, és mely elemek kerülhetnek későbbi fejlesztési szakaszba.

A PRD elkészítése ideális esetben nem egyetlen személy feladata, hanem több szerepkör közös munkájának eredménye. A termékoldali szereplők a felhasználói igényeket és az üzleti célokat képviselik, a fejlesztők a technikai megvalósíthatóságot és a becsült ráfordítást vizsgálják, a tervezői és minőségbiztosítási szempontok pedig a használhatóságot, tesztelhetőséget és a követelmények egyértelműségét erősítik. Ez azért fontos, mert egy jól elkészített PRD nemcsak azt rögzíti, hogy mit kell megvalósítani, hanem azt is, hogy a csapat miért éppen ezeket a funkciókat választotta.

A WatchDB esetében a PRD különösen fontos szerepet tölt be, mert a kezdeti termékvízió szélesebb, mint az elsőként megvalósítható alkalmazás. A dokumentum alapján a projekt nemcsak személyes filmnaplóként értelmezhető, hanem hosszabb távon közösségi filmkövető platformként is. Ez a vízió több irányt tartalmaz: a felhasználó saját filmjeinek naplózását, a watchlist kezelését, értékelések rögzítését, valamint későbbi verziókban a barátok aktivitásának követését és a közösségi ajánlások megjelenítését.

A PRD egyik legfontosabb feladata ebben a helyzetben a kezdeti termék terjedelmének a kezelése. Egy teljes termékvízió gyakran több olyan funkciót tartalmaz, amelyek egyszerre történő megvalósítása túl nagy fejlesztési terhet jelentene. Ezért a dokumentum a funkciókat prioritás szerint csoportosítja: az MVP-hez szükséges funkciókra, a későbbi verziókban megvalósítható elemekre, valamint azokra a részekre, amelyek tudatosan kívül maradnak a projekt hatókörén. Ez a felosztás segít elkerülni, hogy az első verzió túl összetetté váljon, miközben megőrzi a termék hosszabb távú fejlesztési irányát.

A WatchDB PRD-je alapján az MVP központi célja, hogy a felhasználó gyorsan és egyszerűen tudjon filmeket keresni, megnézett filmeket naplózni, értékelést adni, valamint külön listában vezetni azokat a filmeket, amelyeket később szeretne megtekinteni. Ezek a funkciók alkotják azt a minimális terméket, amely már önmagában is értéket ad a felhasználónak, ez az MVP (Minimum Viable Product) lényege. A közösségi funkciók, például a publikus profilok, a követési rendszer vagy az aktivitási feed, a termékvízió fontos részei, de az első megvalósítható verzió szempontjából későbbi fejlesztési szakaszba sorolhatóak.

A PRD további előnye, hogy a fejlesztési döntéseket később visszakereshetővé és indokolhatóvá teszi. A célközönség, a felhasználói folyamatok, a nem funkcionális követelmények, a technológiai függőségek és a korlátok mind olyan tényezők, amelyek befolyásolják az alkalmazás tervezését. Például a külső filmadatok használata indokolja a TMDB API integrációját, míg az email-alapú bejelentkezés az MVP egyszerűsítését szolgálja.

Összességében a PRD a WatchDB fejlesztésében hidat képez a termékötlet és a technikai megvalósítás között. Segítségével a projekt nem elszigetelt funkciók gyűjteményeként jelenik meg, hanem olyan tudatosan szűkített fejlesztési feladatként, amely egy nagyobb termékvízió első, megvalósítható lépését képviseli és egy olyan irányt mutató dokumentum, amihez a csapat könnyen tud igazodni és utalni rá a fejlesztési folyamat során.

\section{Termékvízió}

A WatchDB termékvíziója egy olyan webalkalmazás létrehozása, amely egyszerre működik személyes filmnaplóként és később közösségi filmfelfedező platformként. A projekt alapgondolata, hogy a felhasználó egyetlen, platformfüggetlen helyen tudja rögzíteni, milyen filmeket nézett meg, miket szeretne később megnézni, és hogyan értékelte a már látott filmeket. Ez a megközelítés elsősorban a filmnézési előzmények rendszerezésére és a személyes filmnapló kialakítására épül.

A termék egyik célja a személyes használat. Ebben a felhasználó filmeket kereshet, részletes információkat tekinthet meg róluk, hozzáadhatja őket a watchlisthez, illetve megnézettként jelölheti és értékelheti őket. Ez a funkciókör adja a termék alapértékét, mert akkor is hasznos, ha a felhasználó egyedül használja az alkalmazást. A cél az, hogy a filmek naplózása gyors és egyszerű legyen, ne igényeljen hosszú szöveges értékeléseket, és ne legyen bonyolultabb annál, mint amit egy hétköznapi filmnéző szívesen használna.

A termék második, hosszabb távú célja a közösségi használat. A PRD alapján a WatchDB nemcsak azt tenné lehetővé, hogy a felhasználó saját filmjeit nyilvántartsa, hanem azt is, hogy lássa, barátai vagy ismerősei milyen filmeket néztek meg, hogyan értékelték azokat, illetve milyen filmeket tettek a saját listáikra. Ez a közösségi dimenzió a személyes filmnaplót ajánlási és felfedezési felületté bővítené: a felhasználó nem algoritmikus ajánlórendszeren keresztül, hanem ismerőseinek aktivitása alapján találhatna új filmeket.

Fontos különbséget tenni a teljes termékvízió és az elsőként megvalósítandó MVP között. A teljes vízió része lehet a publikus profil, a követési rendszer, az aktivitási feed, a reakciók, a statisztikai nézetek vagy akár egy fejlettebb ajánlórendszer is. Ezek azonban nem feltétlenül szükségesek ahhoz, hogy az alkalmazás első verziója értéket adjon. Az MVP ezért a termék személyes filmnapló funkcióira koncentrál: a keresésre, a filmrészletek megjelenítésére, a watchlist kezelésére, a megnézett filmek naplózására és az értékelésre.
\section{Célközönség és perszónák}

A célközönség meghatározása a termékkoncepció egyik fontos része, mert a fejlesztési döntések nem választhatók el attól, hogy kiknek készül az alkalmazás. A PRD alapján a termék elsődleges célcsoportját azok az alkalmi filmnézők alkotják, akik havonta néhány filmet néznek, és egyszerű, gyors megoldást keresnek a filmjeik nyilvántartására. Másodlagos célcsoportként megjelennek azok a filmrajongók is, akik tudatosabban követik filmnézési szokásaikat, és hosszabb távon részletesebb statisztikákat vagy közösségi funkciókat is használnának.

Az első célszemély az alkalmi filmkövető típusa. Ő rendszeresen néz filmeket streaming szolgáltatásokon vagy moziban, de nem feltétlenül tekinti magát filmrajongónak. Számára a legfontosabb érték az egyszerűség, gyorsan szeretné rögzíteni, hogy mit látott már, mit szeretne később megnézni, és esetleg milyen benyomása volt egy adott filmről. Ez a felhasználói típus indokolja, hogy az MVP ne legyen túlterhelt bonyolult funkciókkal, hosszú értékelési lehetőségekkel vagy összetett közösségi funkciókkal. A keresésnek, a megnézendő filmek listájának kezelésének és a megnézett filmek naplózásának ezért kevés lépésből kell állnia.

A második célszemély a közösségi filmnéző, aki nemcsak saját magának vezetne listát, hanem kíváncsi arra is, hogy ismerősei milyen filmeket néznek és hogyan értékelik azokat. Ennél a felhasználói típusnál a filmnézés egy bizonyos szociális aspektussá is válik. A PRD-ben szereplő későbbi közösségi funkciók, például a publikus profil, a követési rendszer és az aktivitási feed ehhez a célcsoporthoz kapcsolódik. Ezek a funkciók nem feltétlenül részei az első MVP-nek, de fontosak a termék számára hosszútávon.

A harmadik célszemély a tudatos filmrajongó vagy filmkedvelő, aki nagyobb mennyiségben néz filmeket, szívesen értékel és rendszerez, de nem feltétlenül szeretne hosszú kritikákat írni. Számára a személyes profil, a megtekintett filmek listája, az értékelések és a későbbi statisztikai kimutatások bizonyulhatnak értékesnek. Ez a perszóna bemutatja, hogy a termékvízióban a filmnapló mellett megjelenhetnek mélyebb funkciók is, például műfaji statisztikák, éves összegzések, egyéni listák vagy ajánlási lehetőségek.

A célszemélyek közös jellemzője, hogy mindhárom esetben fontos a gyors használhatóság és az átlátható felület. A különbség inkább abban jelenik meg, hogy az egyes felhasználói típusok milyen mélységben használnák az alkalmazást. Az alkalmi felhasználó számára az MVP alapfunkciói is elegendő értéket adhatnak, míg a közösségi filmnéző és a filmrajongó inkább a későbbi verziókban megjelenő bővítéseket tartaná hasznosnak. Ez a felosztás segít abban, hogy az első fejlesztési szakasz a legszélesebb célcsoport számára hasznos funkciókra koncentráljon, miközben a termék jövőbeli bővíthetőségére is lehetőség adódik.

Ennek alapján a WatchDB célközönsége nem egyetlen, szűken meghatározott, hanem több, eltérő igényű és filmnézési szokásokkal rendelkező csoportból áll. A PRD-ben bemutatott célszemélyek szerepe ezért az, hogy a funkciók fontossági sorrendje először a legnagyobb célközönség alapvető igényeihez igazodjon, és csak ezt követően kerüljenek előtérbe a szűkebb felhasználói csoportokat kiszolgáló bővítések.

\section{Piaci kontextus}

A WatchDB termékkoncepciójának értelmezéséhez érdemes röviden áttekinteni azokat a meglévő megoldásokat, amelyek részben hasonló problémára adnak választ. A piaci kontextus célja ebben az esetben nem egy teljes versenytárselemzés, hanem annak bemutatása, hogy a filmek követése, listázása és megosztása már létező felhasználói igény, ugyanakkor a különböző szolgáltatások eltérő megoldásokkal közelítik meg ezt a problémát.

Az IMDb egyik fontos releváns funkciója a watchlist. A felhasználók címeket adhatnak a saját watchlistjükhöz, rendezhetik azt, illetve értékelések és listák alapján is fedezhetnek fel új tartalmakat \cite{imdbWatchlistFaq}. Ez a megoldás jól mutatja, hogy a filmek későbbi megtekintésre való elmentése alapvető igény. A WatchDB szempontjából ugyanakkor az IMDb inkább információközpontú példaként jelenik meg, míg a projekt célja egy letisztultabb, személyes filmnapló és később közösségi élmény köré szerveződő alkalmazás.

A Trakt erősebb követési és integrációs szemléletet képvisel. A szolgáltatás filmek és sorozatok követésére, listák kezelésére, ajánlásokra, valamint különböző alkalmazásokkal és médiaközpontokkal való összekapcsolásra épít \cite{traktDiscoverTrackShare}. Ez a megközelítés funkcionálisan erős, ugyanakkor összetettebb termékélményt eredményezhet. A WatchDB PRD-je ezért inkább egy egyszerűbb, filmközpontú induló verziót határoz meg, ahol az elsődleges cél nem a sok integráció, hanem a gyors filmkeresés, naplózás és listaépítés.

A piacon jelen vannak kimondottan közösségi filmnapló platformok is, például a Letterboxd. A Letterboxd a filmek nyilvántartását, értékelését, naplózását, listázását és a barátok követését is támogatja, vagyis erősen kapcsolódik ahhoz a közösségi irányhoz, amely a WatchDB hosszabb távú víziójában is megjelenik \cite{letterboxdSocialFilmDiscovery}. Ez azt mutatja, hogy a filmnapló és a közösségi filmfelfedezés kombinációja valós felhasználói igényre épül. A WatchDB dolgozatbeli szerepe azonban nem az, hogy minden létező hasonló szolgáltatással teljes funkcionalitásban versenyezzen, hanem az, hogy egy ilyen termék tervezési és fejlesztési folyamatát egy szűkített MVP-n keresztül mutassa be.

A PRD-ben szereplő piaci összehasonlítás alapján a WatchDB lehetősége abban áll, hogy a fenti irányokból egy tudatosan szűkített változatot határoz meg. Az IMDb adatgazdagságából, a Trakt követési szemléletéből, a streaming szolgáltatások kényelmes listakezeléséből és a Letterboxd közösségi megközelítéséből is levonhatók tanulságok. Az MVP szempontjából azonban ezek közül azok a funkciók kapnak elsőbbséget, amelyek a legszélesebb célközönség számára azonnali értéket adnak: filmkeresés, watchlist, megnézett filmek naplózása és egyszerű értékelés. A közösségi és haladó funkciók így nem elvetett elemek, hanem a termék későbbi fejlődési irányai.

\section{Termékcélok}

A termékcélok meghatározása a PRD egyik kulcsfontosságú része, mert ezek jelölik ki, hogy a fejlesztés során milyen értéket kell létrehozni a felhasználók számára. A célok nem azonosak a konkrét funkciókkal: inkább azt fogalmazzák meg, hogy az alkalmazásnak milyen problémára kell megoldást adnia, milyen felhasználói élményt kell támogatnia, és milyen szempontok alapján lehet eldönteni, hogy az első verzió sikeresnek tekinthető-e.

A WatchDB elsődleges termékcélja, hogy a felhasználó gyorsan és egyszerűen tudjon filmeket keresni, majd azokat saját listáihoz hozzáadni. Ez a cél közvetlenül kapcsolódik a projekt kiinduló problémájához: a filmnézési élmények több platformon és helyzetben történnek, ezért a felhasználónak szüksége van egy egységes, platformfüggetlen nyilvántartásra. Az MVP-ben ennek a célnak a filmkeresés, a filmrészletek megjelenítése, a watchlist és a megnézett filmek naplózása felel meg.

A második fontos cél a személyes filmnapló felépítése. A WatchDB nem csupán egy megnézendő filmeket tartalmazó lista, hanem olyan személyes filmnapló, amelyben a felhasználó visszakeresheti, milyen filmeket látott, hogyan értékelte őket, és milyen filmeket szeretne később megnézni. Ez a cél azért fontos, mert az alkalmazásnak akkor is értéket kell adnia, ha a felhasználó egyelőre nem használ közösségi funkciókat. Az MVP tehát önmagában is használható termékként jelenik meg, nem csak egy későbbi, nagyobb rendszer előzetes változataként.

A harmadik termékcél a közösségi filmfelfedezés előkészítése. A teljes termékvízió alapján a WatchDB hosszabb távon lehetőséget adna arra, hogy a felhasználók kövessék egymás aktivitását, megtekintsék barátaik értékeléseit, és ezek alapján fedezzenek fel új filmeket. Ez a cél az MVP-ben még nem feltétlenül jelenik meg teljes funkcionalitásként, de fontos tervezési szempontként hat a rendszerre. A felhasználói profil, a naplózott filmek strukturált tárolása és a később bővíthető adatmodell mind olyan elemek, amelyek támogatják ezt az irányt.

A negyedik cél a gyors, reszponzív és könnyen használható felhasználói élmény biztosítása. Mivel a célközönség jelentős része alkalmi filmnéző, az alkalmazásnak nem szabad túl sok döntést vagy bonyolult lépést kérnie a felhasználótól. A film keresése, listához adása vagy értékelése akkor működik jól, ha néhány kattintással elvégezhető. Ez a cél nemcsak felületi kérdés, hanem technikai döntéseket is befolyásol: fontos a gyors oldalbetöltés, a hatékony keresés, a mobilon is használható felület és az egyértelmű navigáció.

A PRD-ben szereplő sikerességi mutatók ezekhez a célokhoz kapcsolódnak. Ilyen mutató lehet például a regisztrált felhasználók száma, az aktív felhasználók aránya, az egy felhasználóra jutó havi naplózott filmek száma, a kereséstől a naplózásig eltelt idő, illetve későbbi verzióban az is, hogy hány felhasználó kapcsolódik legalább egy ismerőshöz. Ezek a mérőszámok azért hasznosak, mert a termékcélokat ellenőrizhetővé teszik: nemcsak azt mutatják meg, hogy elkészültek-e a funkciók, hanem azt is, hogy a felhasználók valóban használják-e őket. Emellett a termékcélok a későbbi iteratív fejlesztésben is támpontot adhatnak: segíthetnek eldönteni, hogy érdemes-e új funkcióval bővíteni az alkalmazást, vagy éppen azt, hogy egy kevéssé használt funkció fenntartása továbbra is indokolt-e.

Összességében a WatchDB termékcéljai a személyes értékteremtésre és a későbbi közösségi bővíthetőségre épülnek. Az első verzióban a hangsúly a gyors filmkeresésen, a watchlist kezelésén, a megnézett filmek naplózásán és az egyszerű értékelésen van. Ezek a célok elég szűkek ahhoz, hogy egy MVP keretében megvalósíthatók legyenek, ugyanakkor elég stabil állapotot adnak ahhoz, hogy a termék később közösségi filmfelfedező platformmá fejlődhessen.

\section{Sikerességi metrikák}

A termékcélok mellett a PRD sikerességi metrikákat is meghatároz. Ezek szerepe, hogy a célokat mérhetőbbé tegyék, és a fejlesztés későbbi szakaszaiban visszajelzést adjanak arról, hogy az alkalmazás valóban betölti-e a tervezett szerepét. Egy funkció elkészülése önmagában még nem jelenti azt, hogy a termék sikeres: azt is vizsgálni kell, hogy a felhasználók megtalálják-e az alkalmazást, rendszeresen visszatérnek-e, és ténylegesen használják-e azokat a funkciókat, amelyekre az MVP épül.

A WatchDB esetében a sikerességi mutatók három fő terület köré rendezhetők. Az első a felhasználói növekedés, amely azt mutatja meg, hogy az alkalmazás képes-e elérni egy kezdeti felhasználói bázist. A második az aktivitás és megtartás, amely azt vizsgálja, hogy a regisztrált felhasználók visszatérnek-e, illetve rendszeresen naplóznak-e filmeket. A harmadik a használati élmény és teljesítmény, amely azt méri, hogy a legfontosabb felhasználói folyamatok, például a filmkeresés és a naplózás, elég gyorsan és gördülékenyen működnek-e.

\begin{table}[H]
\centering
\renewcommand{\arraystretch}{1.25}
\begin{tabular}{|p{0.32\textwidth}|p{0.26\textwidth}|p{0.26\textwidth}|}
\hline
\textbf{Metrika} & \textbf{Cél 3 hónap után} & \textbf{Cél 6 hónap után} \\
\hline
Regisztrált felhasználók száma & 500 felhasználó & 2000 felhasználó \\
\hline
Aktív felhasználónként naplózott filmek száma havonta & Legalább 5 film & Legalább 8 film \\
\hline
Heti aktív felhasználók aránya & A regisztrált felhasználók 30\%-a & A regisztrált felhasználók 40\%-a \\
\hline
Legalább egy ismerőssel rendelkező felhasználók aránya & Nem elsődleges MVP-metrika & A felhasználók 60\%-a \\
\hline
Kereséstől naplózásig eltelt átlagos idő & $<$ 15 másodperc & $<$ 10 másodperc \\
\hline
Oldalbetöltési idő & $<$ 2 másodperc & $<$ 2 másodperc \\
\hline
\end{tabular}
\caption{A WatchDB PRD-ben meghatározott sikerességi metrikái}
\label{tab:watchdb-success-metrics}
\end{table}

A regisztrált felhasználók száma elsősorban azt jelzi, hogy a termék képes-e érdeklődést kiváltani a célközönségből. Ez a metrika azonban önmagában nem elegendő, mert egy alkalmazás használhatóságát és valódi értékét nem pusztán a regisztrációk száma mutatja meg. Emiatt fontosabb visszajelzést adhat az aktív felhasználók aránya és a naplózott filmek száma. Ha a felhasználók havonta több filmet is rögzítenek, az arra utal, hogy az alkalmazás valóban beépülhet a filmnézési szokásaikba.

A kereséstől a naplózásig eltelt idő különösen fontos metrika az MVP szempontjából. A WatchDB egyik alapcélja, hogy a filmek rögzítése gyors és kevés lépésből álló folyamat legyen. Ha ez az idő túl hosszú, az arra utalhat, hogy a keresési felület, a filmrészletek oldala vagy a naplózási művelet túl bonyolult. Ez közvetlenül visszahat a felhasználói élményre, és indokolhatja a felület vagy a felhasználói folyamat egyszerűsítését.

A teljesítményhez kapcsolódó metrikák, például az oldalbetöltési idő, azért fontosak, mert a célközönség jelentős része hétköznapi, gyors használatra alkalmas alkalmazást vár. Egy film keresése vagy listához adása gyakran rövid, spontán művelet, ezért a lassú betöltés vagy akadozó felület könnyen csökkentheti a használati hajlandóságot. A két másodperc alatti oldalbetöltési cél ezért nem csupán technikai elvárás, hanem a termék használhatóságához kapcsolódó követelmény is.

A közösségi metrikák, például az ismerősökhöz kapcsolódó felhasználók aránya, a PRD-ben inkább a későbbi termékverzióhoz kötődnek. Az MVP elsődleges célja még a személyes filmnapló működőképességének bizonyítása, ezért a közösségi kapcsolatok aránya az első három hónapban nem tekinthető alapvető sikerességi feltételnek. Hat hónapos távon viszont már fontos visszajelzést adhat arról, hogy a WatchDB képes-e túllépni az egyéni használaton, és valóban közösségi filmfelfedező platform irányába fejlődni.

Összességében a sikerességi metrikák nemcsak utólagos értékelésre szolgálnak, hanem a fejlesztési döntéseket is támogatják. Segítségükkel mérhetővé válik, hogy az MVP mely részei működnek jól, mely folyamatokon kell egyszerűsíteni, és mely későbbi funkciók bevezetése indokolt. Ez különösen fontos egy olyan termék esetében, amelynek hosszabb távú víziója szélesebb, mint az első megvalósított verzió.

\section{MVP és verziókra bontás}

A PRD egyik legfontosabb döntése az MVP, vagyis a Minimum Viable Product meghatározása. Ennek szerepe, hogy a teljes termékvízióból kiválassza azt a legkisebb, de már önmagában is értéket adó funkciókört, amely alapján az alkalmazás első verziója elkészíthető és kipróbálható. A WatchDB esetében ez különösen fontos, mert a hosszabb távú termékvízió személyes filmnaplót, közösségi aktivitást, ajánlásokat, statisztikákat és többféle listakezelési lehetőséget is tartalmaz. Ha ezek egyszerre kerülnének be az első fejlesztési szakaszba, az jelentősen növelné a projekt komplexitását.

Az MVP célja ezért nem az, hogy a WatchDB teljes vízióját azonnal megvalósítsa, hanem az, hogy bizonyítsa az alapötlet működőképességét. Az első verzióban a felhasználónak képesnek kell lennie regisztrálni, filmeket keresni, watchlistet vezetni, megnézett filmeket naplózni, értékelést adni, valamint saját profilján áttekinteni a rögzített filmeket. Ezek a funkciók közvetlenül kapcsolódnak a projekt kiinduló problémájához: a felhasználó egyetlen helyen szeretné követni, mit látott már, mit szeretne később megnézni, és hogyan értékelte a filmeket.

\begin{table}[H]
\centering
\renewcommand{\arraystretch}{1.25}
\begin{tabular}{|p{0.18\textwidth}|p{0.36\textwidth}|p{0.36\textwidth}|}
\hline
\textbf{Verzió} & \textbf{Fő funkciók} & \textbf{Termékoldali szerep} \\
\hline
MVP & Regisztráció és bejelentkezés, filmkeresés, watchlist, megnézett filmek naplózása, értékelés, saját profil & Az alapérték bizonyítása: a felhasználó önállóan is tud személyes filmnaplót vezetni. \\
\hline
v1.1 & Publikus profilok, követési rendszer, aktivitási feed, reakciók, filmnapló idővonal & A közösségi réteg megalapozása, amely a személyes filmnaplót filmfelfedező felületté bővíti. \\
\hline
v1.2 & Többszöri megtekintés kezelése, kommentek, egyéni listák, statisztikai nézetek, ajánlások, keresési szűrők, OAuth bejelentkezés & Mélyebb felhasználói élmény és nagyobb elköteleződés támogatása a már működő alaprendszerre építve. \\
\hline
Hatókörön kívül & Sorozatkövetés, hosszú szöveges kritikák, streaming-elérhetőség, belső üzenetküldés, natív mobilalkalmazás, moderációs eszközök & A projekt fókuszának megőrzése, hogy az első verzió ne térjen el a személyes filmnapló és későbbi közösségi filmfelfedezés fő irányától. \\
\hline
\end{tabular}
\caption{A WatchDB funkcióinak verziókra bontása}
\label{tab:watchdb-versioning}
\end{table}

Az MVP-be kerülő funkciók kiválasztásánál az elsődleges szempont az volt, hogy az alkalmazás közösségi funkciók nélkül is használható legyen. Ez azért fontos, mert egy közösségi termék akkor működik jól, ha már van elegendő felhasználó és tartalom, amelyre a közösségi réteg épülhet. A WatchDB első verziója ezért először a személyes használati értéket teremti meg: a felhasználó saját filmjeit tudja rögzíteni, listázni és értékelni. Ez stabil alapot ad ahhoz, hogy később a barátok aktivitása, a publikus profilok és az ajánlási lehetőségek is értelmezhetővé váljanak.

A v1.1 verzió a PRD alapján a közösségi alapréteg kialakítására fókuszál. Ebben a szakaszban jelenhetnek meg a publikus profilok, a követési rendszer, az aktivitási feed és a reakciók. Ezek a funkciók már nem csupán a személyes nyilvántartást támogatják, hanem azt is lehetővé teszik, hogy a felhasználó mások filmnézési tevékenységéből inspirációt merítsen. Ez a lépés kapcsolja össze a termék személyes filmnapló jellegét a hosszabb távú közösségi filmfelfedező vízióval.

A v1.2 verzió a mélyebb használatot és az elköteleződést támogató funkciókat tartalmazza. Ide sorolhatók a többszöri megtekintések kezelése, a kommentek, az egyéni listák, a statisztikai nézetek, az ajánlások, a fejlettebb keresési szűrők és az OAuth alapú bejelentkezés. Ezek a funkciók hasznosak lehetnek, de nem szükségesek ahhoz, hogy az első verzió teljesítse az alapfeladatát. Emiatt későbbi fejlesztési szakaszba kerülnek, amikor az MVP működéséből és a felhasználói visszajelzésekből már pontosabban látható, mely bővítések indokoltak.

A hatókörön kívülre sorolt elemek legalább annyira fontosak, mint a megvalósítandó funkciók. A sorozatkövetés, a hosszú kritikák írása, a streaming-elérhetőség megjelenítése, a belső üzenetküldés vagy a natív mobilalkalmazás mind olyan irányok, amelyek önmagukban értékesek lehetnének, de eltérítenék a projektet az elsődleges céltól. Ezek kizárása segít abban, hogy a fejlesztés ne váljon túl szélessé, és az első verzió valóban a filmkeresésre, a watchlistre, a naplózásra és az értékelésre koncentráljon.

Ez a verziókra bontás a dolgozat szempontjából is lényeges, mert jól mutatja, hogy a projekt nem egyszerűen funkciók felsorolásaként épül fel, hanem tudatos terméktervezési döntések eredményeként. Az MVP kijelöli a megvalósítás első határát, a későbbi verziók pedig megőrzik a termék fejlődési irányát. Így a WatchDB nem lezárt végtermékként, hanem egy nagyobb termékvízió első, megvalósítható és értékelhető lépéseként jelenik meg.

\section{A user flow-k szerepe}

A user flow-k, vagyis felhasználói folyamatok szerepe az, hogy a PRD-ben meghatározott funkciókat konkrét használati helyzetekhez kapcsolják. Egy funkció önmagában csak azt írja le, hogy mire képes az alkalmazás, a user flow viszont azt mutatja meg, hogy a felhasználó milyen lépésekben jut el egy cél eléréséig. Ez különösen fontos az MVP tervezésekor, mert segít eldönteni, hogy mely képernyők, útvonalak és műveletek szükségesek ahhoz, hogy az első verzió valóban használható legyen.

A WatchDB esetében a user flow-k a termék legfontosabb használati eseteit modellezik. Az első folyamat egy új felhasználó első filmnaplózását mutatja be: regisztráció, keresés, film kiválasztása, megnézettként jelölés, értékelés, majd a profiloldalon való megjelenés. Ez a folyamat azért kiemelt jelentőségű, mert végigvezeti a felhasználót azon az útvonalon, amely az alkalmazás alapértékét adja. Ha ez a folyamat túl hosszú, bizonytalan vagy nehezen követhető, akkor az MVP egyik legfontosabb célja sérül.

A második folyamat a watchlist használatára épül. Ebben a felhasználó bejelentkezés után filmet keres, hozzáadja azt a megnézendő filmek listájához, majd a profiloldalon ellenőrzi, hogy a film megjelent-e a listában. Ez a folyamat egyszerűnek tűnhet, de terméktervezési szempontból fontos, mert több alapvető funkciót összekapcsol: az autentikációt, a keresést, a filmadatok kezelését, a listaállapot módosítását és a profilnézetet.

\begin{table}[H]
\centering
\renewcommand{\arraystretch}{1.25}
\begin{tabular}{|p{0.22\textwidth}|p{0.42\textwidth}|p{0.26\textwidth}|}
\hline
\textbf{User flow} & \textbf{Fő cél} & \textbf{Kapcsolódó verzió} \\
\hline
Első film naplózása & A felhasználó regisztráció után filmet keres, megnézettként jelöli és értékeli azt. & MVP \\
\hline
Film hozzáadása watchlisthez & A felhasználó későbbi megtekintésre elment egy filmet, majd ellenőrzi azt a profilján. & MVP \\
\hline
Film felfedezése ismerős aktivitásából & A felhasználó a feedben látott ismerősi aktivitás alapján ment el egy filmet. & v1.1 \\
\hline
Ismerős követése & A felhasználó megkeres egy másik profilt, követi azt, majd annak aktivitása megjelenik a feedben. & v1.1 \\
\hline
Egyéni lista megosztása & A felhasználó saját listát készít, publikussá teszi, majd megosztja másokkal. & v1.2 \\
\hline
\end{tabular}
\caption{A WatchDB PRD-ben szereplő fő felhasználói folyamatok}
\label{tab:watchdb-user-flows}
\end{table}

A PRD-ben szereplő user flow-k jól elkülönítik az MVP és a későbbi verziók céljait. Az első két folyamat közvetlenül az MVP-hez kapcsolódik, mert a személyes filmnapló alapműködését írja le. Ezek alapján meghatározható, hogy az első verzióhoz szükség van regisztrációs és bejelentkezési felületre, keresőoldalra, filmrészletek megjelenítésére, watchlist műveletekre, naplózási és értékelési lehetőségre, valamint egy saját profiloldalra. Így a user flow-k nemcsak a felhasználói élményt írják le, hanem közvetlenül segítik az oldaltérkép és a fejlesztési feladatok kialakítását is.

A harmadik és negyedik folyamat már a közösségi irányhoz kötődik. Az aktivitási feedből történő filmfelfedezés és az ismerős követése azt mutatja meg, hogyan válhat a WatchDB személyes filmnaplóból közösségi filmfelfedező platformmá. Ezek a folyamatok fontosak a termékvízió szempontjából, de az MVP hatókörén kívülre kerülnek, mert működésükhöz már publikus profilokra, követési rendszerre, aktivitási adatokra és feed logikára is szükség van.

Az ötödik folyamat, az egyéni lista megosztása, még későbbi bővítési irányt képvisel. Ez már nem az alapvető filmnaplózást, hanem a felhasználói tartalom szervezését és megosztását támogatja. A PRD-ben ezért a v1.2 verzióhoz kapcsolódik, ahol az alkalmazás már nemcsak követésre és naplózásra, hanem mélyebb rendszerezésre és közösségi megosztásra is alkalmas lehet.

A user flow-k fejlesztési szempontból azért is hasznosak, mert ellenőrizhetővé teszik a funkciók működését. Egy-egy folyamat alapján könnyebben megfogalmazhatók az elfogadási feltételek: például sikeres-e a regisztráció után a keresés, megjelenik-e a naplózott film a profiloldalon, vagy a watchlisthez adott film valóban visszakereshető-e. Ez a minőségbiztosítás szempontjából is fontos, mert a tesztelés nem elszigetelt komponensekre, hanem valós felhasználói útvonalakra épülhet.

Összességében a user flow-k a WatchDB tervezésében hidat képeznek a termékkoncepció és az implementáció között. Segítenek abban, hogy a fejlesztés ne pusztán funkciók megvalósítására koncentráljon, hanem arra, hogy a felhasználó ténylegesen el tudja érni a számára fontos célokat. Emiatt a user flow-k nem kiegészítő elemei a PRD-nek, hanem a döntések egyik gyakorlati alapját adják.

\section{Korlátok és függőségek}

A PRD-ben szereplő korlátok és függőségek azt mutatják meg, hogy a terméktervezés nem kizárólag a kívánt funkciókról szól, hanem azokról a feltételekről is, amelyek között a projekt megvalósítható. Egy webalkalmazás fejlesztése során a technológiai, időbeli, pénzügyi és csapatméretből adódó korlátok közvetlenül hatnak arra, hogy mi kerülhet be az első verzióba, milyen architektúra választható, és mely kockázatokra kell előre felkészülni.

A WatchDB esetében az egyik legfontosabb korlát a szűkített MVP-fókusz. A projekt célja nem egy minden funkciót tartalmazó, teljes közösségi platform azonnali megvalósítása, hanem egy olyan első verzió elkészítése, amely bizonyítja a személyes filmnapló alapértékét. Ez a döntés összefügg az időbeli és erőforrásbeli korlátokkal is: a PRD az MVP megvalósítását 6-8 hetes időkeretben képzeli el, kisebb fejlesztői csapattal és alacsony költségvetéssel. Emiatt a fejlesztésnek elsősorban azokra a funkciókra kell koncentrálnia, amelyek nélkül az alkalmazás nem adna önálló értéket.

\begin{table}[H]
\centering
\renewcommand{\arraystretch}{1.25}
\begin{tabular}{|p{0.22\textwidth}|p{0.42\textwidth}|p{0.26\textwidth}|}
\hline
\textbf{Korlát} & \textbf{Jelentése a projektben} & \textbf{Tervezési következmény} \\
\hline
API-korlátok & A külső filmadatok a TMDB API-ból származnak, amelynek elérhetősége és használati korlátai befolyásolják a keresést. & Debounce-olt keresés, gyorsítótárazás és külön adatlekérő réteg alkalmazása indokolt. \\
\hline
Költségvetés & A PRD alacsony költségű, elsősorban ingyenes szolgáltatási csomagokra épülő megvalósítással számol. & Olyan technológiák és szolgáltatások választása szükséges, amelyek kis felhasználószám mellett is fenntarthatók. \\
\hline
Csapatméret & A projekt kisebb, 2-6 fős csapattal számol. & A funkciókörnek szűknek, a fejlesztési folyamatnak pedig jól priorizáltnak kell maradnia. \\
\hline
Időkeret & Az MVP tervezett elkészítése 6-8 hét, a közösségi funkciók pedig későbbi szakaszba kerülnek. & A közösségi és haladó funkciók nem részei az első verziónak, hanem külön iterációkban jelenhetnek meg. \\
\hline
\end{tabular}
\caption{A WatchDB fő projektkorlátai}
\label{tab:watchdb-constraints}
\end{table}

A külső szolgáltatásokhoz kapcsolódó függőségek külön figyelmet igényelnek. A WatchDB működésének egyik alapja a TMDB API, amely a filmadatok, poszterek, címek és leírások forrásaként jelenik meg. Ez gyorsítja a fejlesztést, mert nem szükséges saját filmadatbázist építeni, ugyanakkor kockázatot is jelent: ha az API elérhetősége változik, lelassul, vagy a használati feltételek módosulnak, az közvetlenül érintheti az alkalmazás keresési és filmrészlet-megjelenítési funkcióit.

A hitelesítés és az adatkezelés szintén külső szolgáltatásokra épülhet. A PRD-ben megjelenő Supabase vagy hasonló autentikációs szolgáltató használata egyszerűsíti a bejelentkezés, regisztráció és munkamenet-kezelés megvalósítását. Ez azonban azzal jár, hogy figyelni kell a szolgáltató ingyenes csomagjának korlátaira, az adatvédelmi beállításokra és arra, hogy a rendszer később is használható maradjon.

\begin{table}[H]
\centering
\renewcommand{\arraystretch}{1.25}
\begin{tabular}{|p{0.24\textwidth}|p{0.32\textwidth}|p{0.34\textwidth}|}
\hline
\textbf{Függőség} & \textbf{Kockázat} & \textbf{Lehetséges kezelés} \\
\hline
TMDB API & Elérhetőségi probléma, használati feltételek változása vagy API-korlátok elérése. & Filmadatok helyi gyorsítótárazása és az API-hívások külön szolgáltatásrétegbe szervezése. \\
\hline
Autentikációs szolgáltató & Ingyenes csomag korlátai, szolgáltatófüggőség vagy konfigurációs hiba. & Használat monitorozása, jogosultságok ellenőrzése és hordozható adatmodell kialakítása. \\
\hline
Hosting szolgáltató & Sávszélességi, teljesítménybeli vagy deployment korlátok. & Képek optimalizálása, lazy loading alkalmazása és a hosting terhelésének figyelése. \\
\hline
\end{tabular}
\caption{A WatchDB külső függőségei és kapcsolódó kockázatai}
\label{tab:watchdb-dependencies}
\end{table}

A korlátok mellett a PRD előfeltevéseket is tartalmaz. Ilyen például, hogy a felhasználók modern böngészőt és stabil internetkapcsolatot használnak, a TMDB továbbra is elérhető marad a projekt számára, az email és jelszó alapú bejelentkezés elegendő az MVP-hez, illetve az alkalmazás kezdetben angol nyelvű marad. Ezek az előfeltevések nem technikai részleteknek tűnnek, mégis fontosak, mert meghatározzák, milyen helyzetekre kell az első verzióban felkészülni, és mely igények kezelhetők későbbi fejlesztési szakaszban.

A nem funkcionális követelmények szintén szorosan kapcsolódnak a korlátokhoz és függőségekhez. A gyors oldalbetöltés, a biztonságos munkamenet-kezelés, a mobilon is használható felület és a skálázhatóság nem különálló elvárások, hanem olyan minőségi szempontok, amelyek a választott technológiákra és külső szolgáltatásokra is hatással vannak. Például a filmkeresés teljesítménye nemcsak a felület kialakításán múlik, hanem az API-hívások kezelésén, a gyorsítótárazáson és az adatlekérések optimalizálásán is.

Összességében a korlátok és függőségek dokumentálása segít abban, hogy a WatchDB fejlesztése reális keretek között maradjon. A PRD ezáltal nemcsak azt határozza meg, hogy mit kell megvalósítani, hanem azt is, hogy milyen feltételek mellett, milyen kockázatokkal és milyen kompromisszumokkal történhet a megvalósítás.

